% ============================================================
% main.tex — IEEE two-column conference paper
% Compile:  pdflatex main  →  bibtex main  →  pdflatex main  →  pdflatex main
% ============================================================
\documentclass[conference]{IEEEtran}
\IEEEoverridecommandlockouts

% --- Citations ---
\usepackage{cite}

% --- Math ---
\usepackage{amsmath,amssymb,amsfonts}

% --- Algorithms ---
\usepackage{algorithm}
\usepackage{algorithmic}

% --- Graphics ---
\usepackage{graphicx}
\graphicspath{{figures/}}          % drop your images in paper/figures/

% --- Tables ---
\usepackage{booktabs}

% --- Code listings ---
\usepackage{listings}
\usepackage{xcolor}

% --- URLs ---
\usepackage{url}

% --- Figure placeholder (remove once real images are added) ---
\newcommand{\figplaceholder}[2][4cm]{%
  \noindent\fbox{%
    \begin{minipage}[c][#1][c]{\dimexpr\columnwidth-2\fboxrule-2\fboxsep\relax}%
      \centering\small\itshape [Figure: #2]%
    \end{minipage}%
  }%
}

% --- Hyperlinks (load last among packages) ---
\usepackage[hidelinks]{hyperref}

% ---- Python listing style ----
\lstdefinestyle{pycode}{
  language        = Python,
  basicstyle      = \ttfamily\scriptsize,
  keywordstyle    = \color{blue!70!black}\bfseries,
  stringstyle     = \color{green!55!black},
  commentstyle    = \color{gray}\itshape,
  numbers         = none,
  backgroundcolor = \color{gray!8},
  frame           = single,
  rulecolor       = \color{gray!40},
  tabsize         = 4,
  captionpos      = b,
  breaklines      = true,
  showstringspaces= false,
  morekeywords    = {self, True, False, None}
}
\lstset{style=pycode}

% ============================================================
\begin{document}

\title{An Open-Source Web-Based Testbed for\\
       Federated Learning Experimentation\\
       with the Flower Framework}

% ------ AUTHOR BLOCK ------
\author{
  \IEEEauthorblockN{Filippos Rafail Papadakis}
  \IEEEauthorblockA{
    \textit{Faculty of Information Technology}\\
    \textit{Brno University of Technology}\\
    Brno, Czech Republic\\
    \href{mailto:xpapadp00@stud.fit.vutbr.cz}{xpapadp00@stud.fit.vutbr.cz}
  }
}
% ---------------------------

\maketitle

% ============================================================
\begin{abstract}
Federated learning lets multiple clients train a shared model without
ever sending their raw data to a central server.
The idea is well-established, but actually running FL experiments is still
harder than it should be: you need to set up an execution backend,
wire server-client communication, implement aggregation logic, and
collect metrics manually before you can test a single algorithmic idea.
This paper describes a web-based testbed built on the Flower framework and
PyTorch that tries to close that gap.
Users configure and launch simulations through a browser, optionally upload
their own algorithm, model, or dataset as a plain Python file, and watch
training metrics update round by round in real time.
A strategy-wrapping mechanism captures per-round accuracy, loss, and
model checkpoints without requiring any changes to user-provided code.
The whole system ships as a Docker container and comes with a GitHub Actions
integration so experiments can be triggered automatically on code push.
A working CIFAR-10 example runs out of the box; switching to a different
strategy like FedProx is a configuration change.
On top of this the testbed adds an AI agent layer: the platform's
capabilities are exposed as a Model Context Protocol server, and a
language model drives them through an approval-gated loop, so an assistant
can inspect results, write a modified strategy, launch the comparison run,
and explain what changed.
Supporting that layer required making non-IID partitioning and per-client
model evaluation first-class, which in turn makes client drift directly
measurable.
The project is released as open source.
\end{abstract}

\begin{IEEEkeywords}
federated learning, testbed, Flower framework, PyTorch, simulation,
distributed machine learning, open source, experiment management,
large language models, AI agents, Model Context Protocol
\end{IEEEkeywords}

% ============================================================
\section{Introduction}
\label{sec:introduction}

A lot of the data generated today sits on devices or institutions that
cannot easily share it.
Hospitals cannot pool patient records across borders; phones accumulate
personal data that users would rather keep local.
Federated learning (FL) was designed for this setting~\cite{mcmahan2017communication}:
instead of collecting raw data centrally, clients train on their own data
and only send model updates to a server that aggregates them.
The server never sees the data itself.

The concept proved useful in practice early on.
Google deployed FL for next-word prediction on Android keyboards,
coordinating millions of devices without storing their typing
history~\cite{bonawitz2019towards}.
Since then the research interest has expanded well beyond mobile
keyboards --- healthcare, finance, and IoT applications are all active
areas~\cite{kairouz2021advances}.

Running FL experiments, however, is still unnecessarily complicated.
Even with a good framework, a researcher needs to write the server loop,
define client training, hook up metrics collection, and manage the
execution backend before they can actually test an idea.
Frameworks like Flower~\cite{beutel2020flower},
TensorFlow Federated~\cite{tensorflow_federated}, PySyft, and FATE have
reduced a lot of that boilerplate, but they are all command-line tools.
There is no standard place to configure an experiment, watch it run,
or compare results visually without writing additional code.

This project fills that gap with a web-based testbed built on top of Flower.
The idea is simple: wrap the Flower simulation API in a full-stack
web application so that the entire workflow --- from uploading a custom
strategy to watching accuracy converge --- happens in a browser.
Custom algorithms, models, and datasets are just Python files the user
uploads; the platform loads them dynamically and runs them without
any changes to the core system.
Everything is stored in a database so experiments are reproducible and
results are accessible after the fact.

The rest of the paper covers background on FL and the frameworks available
(Section~\ref{sec:background}), the system design and architecture
(Section~\ref{sec:design}), the more interesting implementation details
(Section~\ref{sec:implementation}), a sample experiment
(Section~\ref{sec:sample}), and a discussion of what works well,
what does not, and where the project could go next
(Sections~\ref{sec:discussion} and~\ref{sec:conclusion}).

% ============================================================
\section{Background and Related Work}
\label{sec:background}

\subsection{Federated Learning}

McMahan et al.~\cite{mcmahan2017communication} introduced FL as a way to
train neural networks across data distributed over many devices.
In the standard horizontal setting, $N$ clients each hold a local
dataset $\mathcal{D}_k$ and the goal is to minimize the globally
weighted loss:
\begin{equation}
  \min_{\mathbf{w}} \; F(\mathbf{w}) = \sum_{k=1}^{N} \frac{n_k}{n} \, F_k(\mathbf{w}),
  \label{eq:fl_obj}
\end{equation}
where $F_k(\mathbf{w}) = \frac{1}{n_k}\sum_{i \in \mathcal{D}_k}
\ell(\mathbf{x}_i, y_i; \mathbf{w})$ is the local loss on client $k$,
$n_k = |\mathcal{D}_k|$, and $n = \sum_k n_k$.

The algorithm they proposed for this --- Federated Averaging, or
FedAvg --- works in communication rounds.
Each round, the server picks a random fraction $C$ of clients, sends
them the current global model, waits for them to run a few epochs of
local SGD, and then averages the returned weights.
Algorithm~\ref{alg:fedavg} describes this more precisely.

\begin{algorithm}[t]
  \caption{FedAvg~\cite{mcmahan2017communication}}
  \label{alg:fedavg}
  \begin{algorithmic}[1]
    \REQUIRE $K$ clients, $T$ rounds, fraction $C$, local epochs $E$,
             learning rate $\eta$
    \STATE Server initializes $\mathbf{w}_0$ randomly
    \FOR{round $t = 1, \ldots, T$}
      \STATE $S_t \leftarrow$ random subset of $\max(\lfloor C \cdot K \rfloor, 1)$
             clients
      \FOR{each client $k \in S_t$ (in parallel)}
        \STATE $\mathbf{w}_{t}^{k} \leftarrow
               \textsc{ClientUpdate}(k,\, \mathbf{w}_{t-1},\, E,\, \eta)$
      \ENDFOR
      \STATE $\mathbf{w}_{t} \leftarrow
             \sum_{k \in S_t} \tfrac{n_k}{n}\, \mathbf{w}_{t}^{k}$
    \ENDFOR
    \RETURN $\mathbf{w}_T$
    \STATE \hspace{-1.4em}\textbf{function} \textsc{ClientUpdate}$(k, \mathbf{w}, E, \eta)$:
    \FOR{each local epoch $e = 1, \ldots, E$ and each mini-batch $\mathcal{B}$}
      \STATE $\mathbf{w} \leftarrow \mathbf{w} - \eta \,
             \nabla_{\mathbf{w}} \ell(\mathbf{w};\, \mathcal{B})$
    \ENDFOR
    \RETURN $\mathbf{w}$
  \end{algorithmic}
\end{algorithm}

\subsection{Known Problems with the Basic Approach}

FedAvg works well when clients hold similar data, but real-world data
is almost never that clean.
Clients tend to have skewed label distributions or entirely different
feature spaces --- what the FL literature calls non-IID data.
Li et al.~\cite{li2020convergence} showed formally that FedAvg does not
converge in general under non-IID conditions; local updates on different
data distributions drift apart, and the averaged global model ends up
worse than it should be.

FedProx~\cite{li2020fedprox} addresses this by adding a proximal term
$\frac{\mu}{2}\|\mathbf{w} - \mathbf{w}_t\|^2$ to each client's local
objective, which limits how far each client can deviate from the current
global model.
Communication is another constraint: for large models, sending full
weight vectors every round is expensive.
Kone\v{c}n\'{y} et al.~\cite{konecny2016strategies} propose structured
and sketched updates to cut bandwidth, while adaptive server
optimizers like FedAdam~\cite{reddi2021adaptive} can reduce the number
of rounds needed to converge.

Privacy is a separate concern.
Even though raw data never leaves the client, gradient updates have
been shown to leak information about training
samples~\cite{kairouz2021advances}.
Differential privacy (DP)~\cite{abadi2016deep} is the standard mitigation:
Gaussian noise is added to updates before aggregation, giving a formal
$(\varepsilon, \delta)$-DP guarantee at the cost of some accuracy.

\subsection{Existing Frameworks}

Several open-source frameworks exist for building FL systems.
Table~\ref{tab:frameworks} gives a quick comparison.
Flower~\cite{beutel2020flower} is the one used in this project.
It is framework-agnostic (works with PyTorch, TensorFlow, JAX, and others),
has a clean strategy interface for implementing custom aggregation logic,
and includes a simulation backend built on Ray~\cite{moritz2018ray} that
can run hundreds of virtual clients on a single machine.
TF Federated~\cite{tensorflow_federated} is Google's production system
but is tightly coupled to TensorFlow.
PySyft focuses on secure computation.
FATE targets enterprise vertical FL scenarios.

None of these provide a web interface for managing experiments.
They are all driven by Python scripts, which is fine for experienced
users but adds friction for anyone who just wants to test an idea quickly
or see results without writing extra code.

\begin{table}[t]
  \centering
  \caption{Comparison of federated learning frameworks}
  \label{tab:frameworks}
  \renewcommand{\arraystretch}{1.25}
  \begin{tabular}{@{}lcccc@{}}
    \toprule
    \textbf{Framework} & \textbf{Sim.} & \textbf{Multi-FW} &
    \textbf{Web UI} & \textbf{License} \\
    \midrule
    Flower~\cite{beutel2020flower}    & $\checkmark$ & $\checkmark$ &
      $\times$         & Apache~2.0 \\
    TF Federated~\cite{tensorflow_federated} & $\checkmark$ & $\times$ &
      $\times$         & Apache~2.0 \\
    PySyft                            & Limited      & $\checkmark$ &
      $\times$         & Apache~2.0 \\
    FATE                              & $\times$     & Limited      &
      Partial          & Apache~2.0 \\
    \midrule
    \textbf{Ours}                     & $\checkmark$ & Ext.         &
      $\checkmark$     & Open-source \\
    \bottomrule
  \end{tabular}
\end{table}

% ============================================================
\section{System Design}
\label{sec:design}

\subsection{Goals}

The design was driven by a few practical requirements.
First, the system should not require users to write any infrastructure
code --- configuring an experiment through a web form and clicking
a button should be enough to launch a simulation.
Second, it should be easy to swap in a custom algorithm, model, or
dataset; the chosen approach is to let users upload plain Python files
that the platform loads at runtime.
Third, it has to show what is happening during a simulation in real time,
since that is the whole point of having a UI.
Fourth, experiment configurations and all outputs (metrics, logs,
checkpoints) should be persisted so results are reproducible and
accessible later.
Finally, the deployment should be simple --- the target was a system
that works with just \texttt{docker compose up}, with no external
dependencies beyond Docker itself.

\subsection{Architecture}

The system is split into three layers, shown in Fig.~\ref{fig:architecture}.

The \textbf{web frontend} is a Next.js/React application.
It communicates with the backend over a standard REST API and receives
live updates during a simulation through Server-Sent Events (SSE).

The \textbf{API and database layer} is handled by Next.js API routes
backed by PostgreSQL.
All experiment state --- configurations, status, metrics, checkpoint
paths --- lives here.
Drizzle ORM is used for type-safe database access.

The \textbf{execution engine} is a Python worker process.
In development it runs as a subprocess of the Next.js server;
in production it runs in an isolated Docker container spawned via the
Docker socket, sharing volumes with the main application.
The worker reads the experiment configuration from the database, loads
any user-uploaded modules, runs the Flower simulation through Ray,
and writes results back.

\begin{figure}[t]
  \centering
  % Replace with: \includegraphics[width=0.97\columnwidth]{architecture}
  \figplaceholder[5cm]{System architecture diagram --- add figures/architecture.png}
  \caption{Three-tier architecture. The frontend talks to the API over
    REST and SSE. The API manages state in PostgreSQL and spawns worker
    containers. Workers run Flower simulations via Ray and write results
    back to the database.}
  \label{fig:architecture}
\end{figure}

\subsection{Web Frontend}

The frontend is built with Next.js~15, React~19, TypeScript, and
TailwindCSS, with Recharts handling the training curve plots.
There are three main pages.

The \textit{experiment creation} page lets users set hyperparameters
(number of clients and rounds, client fraction, local epochs, learning
rate, CPU/GPU allocation per client) and upload up to four Python files:
one each for the algorithm, model, dataset, and configuration.
Template files for each are available to download directly from the page.

The \textit{experiment monitoring} page shows the current round,
live accuracy and loss charts, a list of saved checkpoints with
individual download links, and a scrollable log.
Fig.~\ref{fig:dashboard} shows this during an active run.

The \textit{resource dashboard} displays available CPUs and GPUs and
lists all recent experiments.

\begin{figure}[t]
  \centering
  % Replace with: \includegraphics[width=0.97\columnwidth]{dashboard}
  \figplaceholder[5cm]{Monitoring dashboard screenshot --- add figures/dashboard.png}
  \caption{Experiment monitoring dashboard during a FedAvg simulation.
    Charts update in real time as each round completes.}
  \label{fig:dashboard}
\end{figure}

Real-time updates use SSE rather than WebSockets.
The frontend keeps a persistent connection open to
\texttt{/api/experiments/[id]/stream}, which polls the database every
two seconds and pushes any new metrics or checkpoint records to the client.
It is simpler to implement than WebSockets and sufficient for this use case,
since two-second granularity is fine for monitoring training rounds that
typically take several seconds each.

\subsection{Database Schema}

The schema has four tables.
\textit{experiments} holds hyperparameters, file paths, and the current
status (\texttt{pending}, \texttt{running}, \texttt{completed}, or
\texttt{failed}).
\textit{metrics} stores per-round accuracy and loss values, including a
JSON column for individual client breakdowns.
\textit{modelCheckpoints} maps each round to the path of the saved
\texttt{.pt} file on disk, along with its accuracy and loss.
\textit{clients} records partition sizes for each virtual client.
All tables use UUID primary keys so records remain portable across
different deployment environments.

% ============================================================
\section{Implementation}
\label{sec:implementation}

\subsection{Flower Integration and Strategy Wrapping}
\label{sec:impl_flower}

One challenge that came up early was how to collect metrics and save
checkpoints without asking users to modify their strategy code.
If the platform required that, uploading a custom strategy would mean
adding boilerplate lines just to make it work with the testbed, which
defeats the purpose of keeping user files simple.

The solution is a decorator that wraps whatever strategy the user
provides.
It overrides \texttt{aggregate\_fit()} and \texttt{aggregate\_evaluate()},
delegates to the inner strategy, and then fires a callback with the
aggregated result.
The callback saves the checkpoint and writes the metrics to the database.
Listing~\ref{lst:wrapper} shows the core of this.

\begin{lstlisting}[style=pycode, caption={Strategy decorator that intercepts
    aggregation results to save metrics and checkpoints.},
    label=lst:wrapper]
class _WrappedStrategy(Strategy):
    def __init__(self, strategy, on_aggregate):
        self._inner = strategy
        self._callback = on_aggregate  # saves metrics + checkpoint

    def aggregate_fit(self, rnd, results, failures):
        aggregated, metrics = self._inner.aggregate_fit(
            rnd, results, failures)
        if aggregated is not None:
            self._callback(rnd, aggregated, metrics)
        return aggregated, metrics

    # aggregate_evaluate is wrapped identically
\end{lstlisting}

The actual simulation runs through \texttt{flwr.simulation.run\_simulation()},
which takes a \texttt{ServerApp} and a \texttt{ClientApp} and hands off
scheduling to Ray.
Ray allocates virtual clients to CPU and GPU cores according to the
\texttt{client\_resources} dictionary, which is populated from the
experiment's configuration.
Device selection (CUDA, Apple MPS, or CPU) happens automatically at
startup.

\subsection{Dynamic Module Loading and Validation}
\label{sec:impl_modules}

Uploaded Python files are imported at runtime using \texttt{importlib}.
Before any import happens, the file is parsed with Python's \texttt{ast}
module to check that the expected entry points exist:
\begin{itemize}
  \item \textbf{Algorithm}: must define \texttt{get\_strategy()}
  \item \textbf{Model}: must define \texttt{get\_model()} or a class
        named \texttt{Net} or \texttt{Model}
  \item \textbf{Dataset}: must define
        \texttt{load\_data(partition\_id, num\_partitions)}
  \item \textbf{Config}: must define a \texttt{CONFIG} dict or
        \texttt{get\_config()}
\end{itemize}
This catches obvious mistakes --- a missing function, a typo in the name
--- before the simulation process even starts, so the user gets an error
immediately rather than after waiting for Ray to initialize.
Validated modules are copied to \texttt{.module\_cache/} so that Ray
worker processes, which run in separate Python interpreters, can import
them by name.

\subsection{Experiment Lifecycle}

When a user submits an experiment, the API creates a database record
with status \texttt{pending} and runs the AST validation on any
uploaded files.
If validation passes, a worker is spawned.
The worker reads the configuration, updates status to \texttt{running},
and starts the simulation.
After each round completes, the wrapped strategy fires the callback,
which writes metrics and a checkpoint path to the database --- this is
what the frontend's SSE stream picks up.
When the simulation finishes, the worker saves the execution logs and
sets the status to \texttt{completed} or \texttt{failed}.

\subsection{Deployment and CI/CD Integration}

The Docker image is built in three stages: install Node dependencies,
build the Python virtual environment and compile the Next.js app,
then assemble a minimal runtime image.
The \texttt{docker-compose.yml} brings up the application and a
PostgreSQL~15 container with named volumes for checkpoints, uploads,
and the Hugging Face dataset cache.

The project also ships a reusable GitHub Action in \texttt{gh-action/}.
Any repository can include it to automatically trigger an experiment
when code is pushed to a configured folder (e.g.,
\texttt{flower-simulation/}).
The action finds the relevant Python files by naming convention and
submits them to the testbed.
Optionally it can wait for the simulation to finish and report the
final accuracy and loss as output variables, which makes it possible
to track model performance across commits in the same CI pipeline.

% ============================================================
\section{Sample Application}
\label{sec:sample}

\subsection{Default Experiment: CIFAR-10 with FedAvg}

To make the testbed usable without uploading anything, it ships with
default implementations for model, dataset, and strategy.
The default model is a small CNN~\cite{paszke2019pytorch} with two
convolutional layers (32 and 64 filters, $3\times3$ kernels, ReLU
and max-pooling), followed by three fully connected layers for the
10 CIFAR-10 classes.
The default dataset loader partitions CIFAR-10 into IID shards using
Flower's \texttt{IidPartitioner} from the \texttt{flwr-datasets}
library, with an 80/20 train--test split per partition.
The default strategy is FedAvg with SGD (momentum~$0.9$, learning
rate~$0.01$, batch size~$32$).

A typical run uses 10 virtual clients, 10 rounds, client fraction
$C = 0.3$, and 3 local epochs per round.
These can all be changed from the web form.
Fig.~\ref{fig:metrics} shows how accuracy and loss evolve across
rounds as reported on the monitoring dashboard.

\begin{figure}[t]
  \centering
  % Replace with: \includegraphics[width=0.97\columnwidth]{metrics}
  \figplaceholder[4cm]{Training accuracy/loss chart --- add figures/metrics.png}
  \caption{Per-round accuracy and loss for the default FedAvg experiment
    on CIFAR-10 with 10 virtual clients and $C = 0.3$.}
  \label{fig:metrics}
\end{figure}

\subsection{Switching Strategy: FedProx}

Listing~\ref{lst:fedprox} shows an algorithm file that replaces FedAvg
with FedProx~\cite{li2020fedprox}.
The proximal term $\frac{\mu}{2}\|\mathbf{w} - \mathbf{w}_t\|^2$ it
adds to each client's objective limits how far local updates drift,
which helps when client data distributions are heterogeneous.
Uploading this file and relaunching the experiment is all it takes ---
no other changes are needed.

\begin{lstlisting}[style=pycode, caption={Algorithm file for FedProx.
    Uploading this replaces FedAvg without any other changes.},
    label=lst:fedprox]
import flwr as fl

def get_strategy():
    return fl.server.strategy.FedProx(
        fraction_fit=0.3,
        min_fit_clients=3,
        min_available_clients=10,
        proximal_mu=0.1,     # regularization strength mu
    )
\end{lstlisting}

The same mechanism works for any other Flower strategy, including
FedAdam and FedYogi~\cite{reddi2021adaptive}, and for fully custom
\texttt{Strategy} subclasses.
Model and dataset files can be swapped in the same way to target
different architectures or datasets without touching the platform.


% ============================================================
\section{AI Agent Layer}
\label{sec:agent}

The system described so far removes the boilerplate from running a
federated learning experiment, but it does not help with the part that
actually takes time: deciding what to run next.
Reading a training curve, forming a hypothesis about why accuracy
plateaued, writing a modified strategy, and launching the comparison run
are all still manual.
This section describes an agent layer that automates that loop, built on
top of the same services the web interface uses.

\subsection{Design}

The agent layer adds three components.
A \emph{tool surface} exposes the testbed's capabilities as a set of
typed operations.
An \emph{orchestration loop} lets a large language model call those
operations in sequence, observing each result before deciding the next
step.
A \emph{memory} records what previous experiments showed, so conclusions
survive across sessions.

A deliberate constraint shaped the rest of the design: the agent must not
reach the platform through its own HTTP API.
Instead, both the web routes and the tool surface call a shared service
layer directly.
This keeps a single implementation of every operation, and means an
agent cannot do anything a user could not.

\subsection{Model Context Protocol Server}

The tool surface is exposed as a Model Context Protocol (MCP)
server~\cite{mcp} mounted inside the Next.js application.
MCP is a JSON-RPC protocol for connecting language models to external
systems, and using it rather than a bespoke interface has a practical
consequence: the same server that the built-in agent talks to can be
attached to any MCP-compatible assistant, so a researcher can query their
experiments from a general-purpose coding assistant without the testbed
knowing anything about it.

The server exposes 39 tools in six groups, corresponding to the
capabilities an agent needs:

\begin{itemize}
  \item \textbf{Training control} --- listing, creating, cloning,
        configuring, starting and stopping experiments, plus a blocking
        \texttt{wait\_for\_round} so an agent can follow a run without
        polling.
  \item \textbf{Results} --- per-round metrics, per-client metrics,
        checkpoint inspection, log retrieval, and a computed digest that
        summarises convergence, detects plateaus and divergence, and
        flags conditions that would otherwise be misread.
  \item \textbf{Node management} --- hyperparameters, data partitioning,
        and per-client overrides for modelling heterogeneous nodes.
  \item \textbf{Aggregation strategy} --- selecting and parameterising a
        built-in strategy, or writing and dry-running a custom one.
  \item \textbf{Files} --- reading templates and uploaded modules,
        writing new modules, and validating them.
  \item \textbf{Memory} --- semantic search over past work.
\end{itemize}

Alongside tools, the server exposes experiments, logs, metric series and
templates as MCP \emph{resources}, which an assistant's user can attach
as context directly, and three \emph{prompts} that appear as
slash-commands.

\subsection{Agent Roles and Orchestration}

Rather than running several independent agents, the system defines roles
as modes of a single loop.
Four separate agents would mean four conversation histories and four
times the token cost for what is conceptually one task.
What genuinely differs between roles is the instruction block, the subset
of tools in reach, and how much deliberation the task warrants, so only
those vary: a \emph{planner} with read-only access, an
\emph{experimenter} holding the training and node tools, an
\emph{analyst} holding results and memory, and an \emph{engineer} holding
the code and strategy tools.

The loop itself is conventional: send the conversation to the model,
receive either a final answer or a set of tool calls, execute them,
append the results, and repeat.
Two properties of the implementation are worth noting.
First, the loop runs detached from the HTTP request that started it, so
closing the browser cannot abort a turn midway through a tool call.
Second, all state needed to resume lives in PostgreSQL rather than in
process memory, which is what makes the approval mechanism described
next possible.

\subsection{Approval and Autonomy}

Each tool is classified as \emph{read}, \emph{write} or \emph{execute}.
Read operations run immediately.
Write and execute operations are recorded as pending and the turn
suspends, surfacing a card in the interface that states the action in
plain language, shows a diff for file writes, and keeps the exact
request one click away.
Approving resumes the turn; declining also resumes it, with a tool
result saying the user declined.
That detail matters more than it appears: a tool call left without a
matching result is a protocol error, so a declined action must still
produce one, or the conversation cannot continue.

Classification is fail-closed --- an unrecognised tool is treated as
\emph{execute} and gated --- and enforced on the server, never by
instructing the model.
A per-conversation toggle lifts the gate for users who want the agent to
act without interruption.

\subsection{Memory}

Completed experiments are summarised and embedded into a
pgvector~\cite{pgvector} column, and relevant entries are retrieved and
injected into the user's message before each turn.
Embedding models differ in output width, so vectors are stored at a fixed
width: narrower models are zero-padded, which is exact for cosine
similarity since padding changes neither dot products nor norms, and
every row records the model that produced it so retrieval filters to the
active one.
Switching embedding model therefore degrades to an empty recall rather
than to meaningless neighbours, and never requires a schema migration.
Because not every provider offers an embeddings endpoint, retrieval falls
back to PostgreSQL full-text search over the same stored text when none
is configured.

\subsection{Capability Changes to the Execution Engine}

Three properties of the original execution engine were adequate for a
human driving the web form but became limitations once an agent was
driving it, and were changed.

The strategy factory returned a user-supplied strategy unmodified,
silently dropping the metric aggregation callbacks the platform attaches
to its own strategies.
A run configured that way trains correctly but records no evaluation
metrics at all, which is indistinguishable from a model that failed to
learn until the run finishes.
The factory now injects the platform defaults for any callback the
strategy leaves unset, and a dry-run tool reports which callbacks a
module defines before the run starts.

The dynamic module loader cached uploaded modules under a fixed name in a
single shared directory.
Two concurrent experiments therefore overwrote each other's modules and
each trained whichever was copied last, without any error.
The cache is now scoped per experiment.
This mattered little when experiments were launched one at a time through
a form; an agent's natural behaviour is to fan out variants.

Finally, the aggregation strategy could previously only be changed by
uploading a Python file.
FedProx, FedAdam, FedAdagrad and FedYogi are now selectable by name with
their parameters exposed, which makes strategy comparison --- the
capability the testbed exists to support --- a configuration change
rather than a programming task.

\subsection{Non-IID Partitioning and Local Model Evaluation}

Two capabilities were added that the agent needs and that the testbed
lacked.

The default dataset loader now supports Dirichlet, shard and pathological
partitioning in addition to IID.
This matters because with an IID split, aggregation strategies are
difficult to tell apart; the differences between them are largely
differences in how they cope with heterogeneous client data.

Per-client metrics are now retained rather than being collapsed into a
weighted mean and discarded, and each client's local model can optionally
be checkpointed alongside the aggregated global one.
A standalone evaluation entry point can then run any saved checkpoint
over any client's data partition.
This makes client drift directly measurable.
In one run with Dirichlet partitioning at $\alpha = 0.3$, a client's
local model reached 61.0\% accuracy on its own partition against the
global model's 57.8\%, while scoring 19.6\% on a different client's
partition where the global model reached 35.2\% --- the local model had
fit its own skewed label distribution at the cost of generalisation.
That comparison was not expressible in the original system.

\subsection{Interface}

The agent is used through a chat page that renders the conversation,
collapses tool calls to a single line each, expands them on demand, and
displays approval cards inline.
When a tool touches an experiment, a live card for that experiment
appears in the transcript, subscribing to the same event stream the
monitoring page uses.
A settings page holds the model connection --- provider, endpoint, API
key and model --- with the key encrypted at rest and never returned to
the browser.
Because the tool surface is provider-neutral, the same agent runs against
a hosted model or a locally served one.

% ============================================================
\section{Discussion}
\label{sec:discussion}

\subsection{What the Testbed is Actually Useful For}

The main benefit compared to using Flower directly is the reduced
setup time.
Getting from an algorithmic idea to a running simulation requires
writing one small Python file and filling out a web form, rather than
setting up a project, wiring up a training loop, and writing metrics
collection code.
The per-round checkpoints and the persistent experiment log also make
it easier to debug runs after the fact, since everything is stored in
the database and remains accessible after the simulation ends.

For students and courses on distributed ML, the web interface lowers
the entry barrier considerably.
It is possible to demonstrate the effect of changing the number of
clients, the client fraction, or the number of local epochs just by
running multiple experiments from the UI without writing any code.

The GitHub Actions integration is useful for more automated workflows.
A team working on a custom FL algorithm can set up the repository so
that a simulation runs automatically on every pull request, with the
final accuracy reported as part of the CI output.

\subsection{How it Compares to Other Tools}

As shown in Table~\ref{tab:frameworks}, the web interface is the main
thing that sets this testbed apart from existing frameworks.
The OARF benchmark suite~\cite{hu2022oarf} addresses a related problem
--- standardized benchmarking of FL algorithms --- but it is a
command-line tool and does not include a monitoring interface.
The testbed is self-hosted and fully open-source, which distinguishes
it from commercial FL platforms that require cloud accounts.

\subsection{Hardware Support and Deployment}

Because Flower's simulation backend uses Ray~\cite{moritz2018ray}
for scheduling, the testbed can scale to a large number of virtual
clients on a single machine by allocating fractional CPU and GPU
resources per client.
The \texttt{cpus\_per\_client} and \texttt{gpu\_fraction\_per\_client}
parameters are exposed directly in the experiment configuration form.
CUDA and Apple MPS are both detected automatically at startup; if
neither is available the system falls back to CPU without any
configuration change.
This makes the testbed usable on a laptop for development and on a
GPU server for longer runs without any code changes.

\subsection{Limitations}

The most significant limitation is that the simulation runs on a
single machine.
Real network latency, packet loss, and heterogeneous device environments
are not modelled.
All virtual clients also share the same Python environment and hardware,
so experiments cannot capture the system heterogeneity that matters in
real deployments.
User-uploaded code runs in the same Python environment as the rest of
the platform, which is fine for trusted users but would need a proper
sandbox before opening the system to arbitrary submissions.
The agent layer sharpens rather than softens this: code the platform
executes may now be written by a language model rather than by the user,
which converts a user-trust assumption into a model-trust assumption.
Several boundaries limit the exposure --- generated files are confined to
a dedicated workspace directory, statically validated before import, and
every write or execution passes through the approval gate --- but none of
them is a sandbox.
Experiment logs are a further consideration, since they are arbitrary
output from user-supplied code that the agent reads while holding tools
that write files and start jobs; content from such sources is labelled as
untrusted in the agent's context, but the effective protection is that
the approval gate and the workspace confinement are enforced in the
server, where no text the model reads can reach them.
Finally, the production worker container shares the application image and
its volume mounts, so it provides isolation of process, not of trust.

% ============================================================
\section{Conclusion}
\label{sec:conclusion}

The goal of this project was to build an environment for testing
federated learning algorithms that is easier to use than running
Flower directly, while still being flexible enough to support custom
strategies, models, and datasets.
The resulting system is a web application that wraps Flower's simulation
API, stores all experiment state in PostgreSQL, and serves a browser-based
dashboard for configuration and monitoring.
A dynamic module loader lets users inject Python code without modifying
the platform; a strategy decorator captures metrics and checkpoints
transparently.
The whole thing deploys with Docker and integrates with GitHub Actions
for automated experimentation.

There are several directions where the system could be extended.
A genuine sandbox for executing uploaded and generated code is now the
most valuable, given that the agent layer makes model-authored code a
routine occurrence rather than an exception.
Differential privacy support~\cite{abadi2016deep} could be added as a
configurable option without major architectural changes.
Supporting additional frameworks beyond PyTorch, such as TensorFlow
Federated or FATE, is architecturally straightforward since the
execution engine is the only framework-specific component.
A true multi-machine mode and side-by-side experiment comparison in the
web interface are further directions worth exploring.
On the agent layer, the most promising extension is letting it propose
and run a sequence of experiments unattended against a stated objective,
which the approval mechanism and memory already make structurally
possible.

Source code, Docker images, and documentation are available at:
\begin{center}
  \url{https://github.com/phrp720/flower-testbed}
\end{center}

% ============================================================
\bibliographystyle{IEEEtran}
\bibliography{references}

\end{document}
